\documentclass[../main.tex]{subfiles}
\begin{document}

% TODO: resolve chapter and section names

\chapter{Mathematical modelling of orographic flows}

In the previous chapter \Cref{chap:FLOWS} the underlying physical laws governing atmospheric flows have been described. In this chapter we describe further \emph{modelling choices} used to obtain numerical solution to our problem.

\section{Initial and boundary conditions}
\label{sec:IBCs}

Let $T>0$ and denote $\Omega \subset \R^2$ to be the set

\begin{equations}[single,!,numberline=all]
\label{eq:OMEGA}
\Omega = \left\{\left(x, y\right) \in \R^2 \middle| x \in \left(-L, L\right), h(x) < y < H \right\},
\end{equations}

where $L > 0$ is the \textcolor{gray}{half-}width of the domain, $H$ is the height of the domain and $h(x)$ is the \emph{mountain profile}. In our simulations, the \emph{Witch of Agnesi} mountain profile \Cref{fig:WITCH-OF-AGNESI} of height $h_m$ and half-width $a$ has been used.

\begin{figure}[htbp]
	\centering
	\begin{tikzpicture}
		\begin{axis}[
			width=0.85\linewidth,
			height=0.5\linewidth,
			axis lines=left,
			xlabel={$x$},
			ylabel={$h(x)$},
			xmin=-6, xmax=6,
			ymin=0, ymax=1.15,
			xtick={-1,0,1},
			xticklabels={$-a$,$0$,$a$},
			ytick={0,0.5,1},
			yticklabels={$0$,$h_m/2$,$h_m$},
			tick label style={font=\footnotesize},
			label style={font=\small},
			clip=false,
			samples=201,
			domain=-6:6,
			]
			\addplot[very thick, cyan!60!black] {1/(x^2+1)};

			\draw[densely dashed, gray!55] (axis cs:-6,1) -- (axis cs:0,1);
			\draw[densely dashed, gray] (axis cs:1,0) -- (axis cs:1,0.5);
			\draw[densely dashed, gray] (axis cs:-1,0) -- (axis cs:-1,0.5);

			\draw[<->, >=Stealth, shorten >=1pt, shorten <=1pt]
			(axis cs:-1,0.5) -- node[fill=white, inner sep=1.5pt, font=\footnotesize] {$2a$} (axis cs:1,0.5);

			\addplot[only marks, mark=*, mark size=1.3pt, cyan!60!black] coordinates {(0,1)};

			\node[font=\footnotesize, anchor=west, gray!55!black] at (axis cs:2.8,0.46) {$h \sim h_m\,(a/x)^2$};
			\draw[->, >=Stealth, gray] (axis cs:3.6,0.34) to[bend right=12] (axis cs:4.3,0.05);
		\end{axis}
	\end{tikzpicture}
	\caption[The Witch of Agnesi mountain profile]{%
		The \emph{Witch of Agnesi} profile $h(x) = h_m a^2/(x^2 + a^2)$ used as the mountain barrier. The maximal height $h_m$ is attained at $x = 0$, while the half-width $a$ fixes the full width at half maximum $2a$. 
	}
	\label{fig:WITCH-OF-AGNESI}
\end{figure}

Denote further

\begin{equations}[columns,!,numberline=all]
\label{eq:BOUNDARY}
\partial \Omega &= \Gamma_L \cup \Gamma_R \cup \Gamma_B \cup \Gamma_T, \\
\Gamma_L &= \left\{-L\right\} \times \left(h(-L), H\right), \\
\Gamma_R &= \left\{L\right\} \times \left(h(L), H\right), \\
\Gamma_B &= \left(-L, L\right) \times \left\{h(x) \right\}, \\
\Gamma_T &= \left(-L, L\right) \times \left\{H \right\}, 
\end{equations}

to be the left, right, bottom and top boundaries. 

We are interested in numerical solutions to the initial-boundary value problem for the Euler system written in the form\sidenote{This form is common in mathematical literature, as the order of the equation and the convective \emph{nonlinearity} is clearly visible.}

\begin{equations}[columns,*]
\partial_t\left(\rho \vb{v}\right) + \div\left(\rho \vb{v} \otimes \vb{v}\right) &= - \grad p + \rho \vb{g}, \, \text{in} \, \left(0,T\right) \times \Omega, \\
\partial_t \rho + \div\left(\rho \vb{v}\right)&= 0, \, \text{in} \, \left(0,T\right) \times \Omega, \\
p &= R \rho T, \\
\vb{v} &= \vb{v}_0, \, \text{on} \, \left\{ 0 \right\} \times \Omega,\\
\vb{v} &= V \vb{e}_x, \, \text{on} \, \left(0,T\right) \times \Gamma_L,\\
\vb{v} \cdot \vb{n} &= 0, \, \text{on} \, \left(0,T\right) \times \Gamma_B,\\
p &= \overline{p}, \, \text{on} \, \left(0,T\right) \times \Gamma_R,\\
&\text{\enquote{radiation condition, on $\left(0,T\right) \times \Gamma_T$}}
\end{equations}

with $\vb{v}_0, \rho_0$ being the initial velocity and density, $\overline{p}$ the hydrostatic profile and $0 < V < c$ the constant \emph{subsonic} inflow velocity magnitude.

% TODO: add initial conditons
% TODO discuss BCs for Euler
% TODO discuss used BCs 
% TODO discuss radiative BC

\subsection{Boundary conditions for the Euler system}
\label{sec:EULER_BC}


\autocite[p.~7-8]{oligerTheoreticalPracticalAspects1978}

\subsection{Radiation boundary condition}
\label{sec:RADIATION}

\subsection{Terrain-following coordinates}
\label{sec:TERRAIN-CORDS}

% TODO: add figures

\begin{figure}[htbp]
	\centering
	\includegraphics[width=0.85\linewidth]{terrain_following}
	\caption[Terrain-following particle lattice]{%
		Seeding of the particle lattice in terrain-following coordinates \autocite{gal-chenUseCoordinateTransformation1975}, with schematic parameters. The reference lattice (a) is mapped by $y^{*} = y + h(x)\,(H - y)/H$ to the terrain-following lattice (b), whose bottom row lies on the mountain surface $h(x)$ and whose coordinate surfaces relax to horizontal at the domain top $H$. The vertical spacing grows as $\exp(Ky/2)$, coarsening the resolution with altitude.
	}
	\label{fig:TERRAIN-FOLLOWING}
\end{figure}

\section{Well-balanced methods}
\label{section:MODELLING}

This chapter provides an introduction to specialized numerical methods for the compressible Euler equations. The motivation for the use of such methods is presented in the first subsection, whereas our implementation is described in \Cref{sec:WELL-BALANCED-SPH}. Results obtained with the use of these methods are then presented in \Cref{chap:EXPERIMENTS}.

\subsection{Preservation of the hydrostatic balance}
\label{sec:INEXACT}

Internal gravity waves can be considered perturbations of the background atmospheric state. It is standard to take the background atmospheric state as \emph{hydrostatic}\sidenote{As we did in the section \Cref{sec:MTN-WVS}.}: $\vb{v} = 0$ and the pressure $p$ and density $\rho$ solve the equation \Cref{eq:HYDROSTATIC-P}

\begin{equations}[single,*]
	\grad p = \rho \vb{g}.
\end{equations}

Suppose now we wish to \emph{numerically} solve the \emph{non-stationary} isentropic compressible Euler system with the initial condition given as the hydrostatic balance, \ie solve\sidenote{We are using a formulation of the Euler equations that is more common in numerical mathematics.}

\begin{equations}[columns,!,numberline=all]
	\partial_t \rho + \div\left(\rho \vb{v}\right) &= 0, \\
	\partial_t\left(\rho \vb{v}\right) + \div\left(\rho \vb{v} \otimes \vb{v}\right) &= - \grad p + \rho \vb{g},\label{eq:NONSTAT-EULER}\\
	\partial_t \theta + \div\left(\theta \vb{v}\right) &= 0, \\
	\vb{v}\left(0, \vb{x}\right) = \vb{0}, \theta\left(0, \vb{x}\right) &= \theta_0\left(\vb{x}\right), \rho\left(0, \vb{x}\right) = \rho_0\exp\left(- \frac{z}{H_s}\right),
\end{equations}

together with given $p = p\left(\theta, \rho\right)$ and some boundary conditions.\sidenote{A rigorous numerical treatment of the boundary conditions is complicated, see \eg \autocite[p.~59]{levequeFiniteVolumeMethods2002}.} \emph{Mathematically} speaking, since at time $t=0$ the system is in hydrostatic equilibrium, the right-hand side of the momentum equation as well as the convective term vanish and the system \emph{remains} in the hydrostatic equilibrium. Yet, in a numerical simulation, generally the hydrostatic equilibrium is \emph{not preserved},

\begin{equations}[single,*]
	\grad_h p_{i} - \rho_i \vb{g} \neq 0,
\end{equations}

where $p_i, \rho_i$ are some approximations of the solution \enquote{at time $i$} obtained by the numerical method and $\grad_{h}$ is some discrete gradient operator. Since the right-hand side of the momentum balance \Cref{eq:NONSTAT-EULER} does not vanish on the \emph{discrete level}, a spurious non-physical force is introduced and the gas moves with \emph{nonzero velocity.}

\subsection{Implications for wave simulations}
\label{sec:WB-WAVES}

The failure to preserve hydrostatic balance is particularly damaging for simulations of internal gravity waves. As established earlier, internal gravity waves are small perturbations of the \emph{background state.} When the hydrostatic balance is inexact, the spurious forces influence the simulation and affect the flow field. 

Käppeli and Mishra demonstrate empirically in \autocite[p.~6]{kappeliWellbalancedFiniteVolume2016a} that, for small perturbations of the velocity field\sidenote{The authors used a sine perturbation of the form \begin{equations}[single,*]
		\delta v = A\sin\left(\omega t\right),
	\end{equations}
with A = \num{e-8}.}
spurious forces dominate the wave field. By contrast, well-balanced schemes give qualitatively correct wave dynamics. For large-amplitude flows, in \autocite[p.~6]{kappeliWellbalancedFiniteVolume2016a} it is shown that the effect is not critical. 
For small-amplitude wave dynamics---relevant to internal gravity waves---exact discrete hydrostatic balance is therefore essential.

\subsection{Well-balanced schemes}
\label{sec:WB-SCHEMES}

Numerical methods that preserve the \emph{analytic} steady solutions on the \emph{discrete} level are called \emph{well-balanced} \autocite{cargoSchemaEquilibreAdapte1994,greenbergWellBalancedSchemeNumerical1996}. The methods had been first used for shallow water equations \autocite{greenbergWellBalancedSchemeNumerical1996} and are nowadays standard in many areas where a numerical solution of \emph{hyperbolic} PDEs is required, \eg atmospheric flows \autocite{kappeliWellbalancedFiniteVolume2016a,chertockWellbalancedNumericalMethod2023,berthonFullyWellbalancedEntropystable2025} and astrophysics \autocite{kappeliWellbalancedMethodsComputational2022}. 

While well-balanced finite-volume and DG schemes are now standard for atmospheric and astrophysical applications, well-balanced SPH formulations remain underdeveloped. This thesis aims to contribute to that gap. In \Cref{sec:WELL-BALANCED-SPH} we derive well-balanced formulations of SPH suitable for isentropic atmospheric flow --- the central contribution of this thesis.


\end{document}


